\documentclass[10pt,twocolumn]{article}

\usepackage[margin=0.72in]{geometry}
\usepackage[T1]{fontenc}
\usepackage[utf8]{inputenc}
\usepackage{amsmath}
\usepackage{amssymb}
\usepackage{booktabs}
\usepackage{array}
\usepackage{tabularx}
\usepackage{xcolor}
\usepackage{listings}
\usepackage{url}
\usepackage[hidelinks]{hyperref}

\newcolumntype{Y}{>{\raggedright\arraybackslash}X}
\newcommand{\ttc}{\mathrm{TTC}}
\newcommand{\vsevenwarning}{\textbf{Outdated for Scientific Recovery V7: this
paper contains no V7 results and must not be updated before a signed aggregate.}}
\Urlmuskip=0mu plus 2mu
\def\UrlBreaks{\do\/\do\-\do\_\do\.}
\emergencystretch=2em
\sloppy
\lstset{
  basicstyle=\ttfamily\scriptsize,
  breaklines=true,
  columns=fullflexible,
  frame=single,
  framerule=0.2pt,
  xleftmargin=0.4em,
  xrightmargin=0.4em
}

\title{\textbf{E-JEPA-TTC: Cache-Free High-Resolution Event Prediction for Auditable TTC Research}}
\author{
KriptaStudios Research Handoff\\
\small Reproducible local report from the \texttt{e-jepa-ttc} repository\\
\small Protocol revision: 2 August 2026
}
\date{}

\begin{document}
\maketitle

\begin{center}
\fcolorbox{red}{white}{\parbox{0.94\linewidth}{\vsevenwarning}}
\end{center}

\begin{abstract}
E-JEPA-TTC investigates time-to-contact estimation from event cameras using
high-resolution dense tokens and multi-horizon joint-embedding prediction. This
revision does not claim state of the art. It combines an exact historical EvTTC
baseline reconstruction, grouped cross-validation, an audit of the public
Garl-TTC/eAP contracts, and a raw on-demand trainer that avoids a projected
455~GiB dense cache. It also adds a dataset-free, falsifiable audit of semantic
shortcut allocation. The historical baseline obtains $0.322892$~s MAE and
$8.1554\%$ relative error on its original validation split. Dense Patch improves
one matched split but loses five-fold, three-seed grouped cross-validation to the
global A0 control. A first real high-resolution 16/16-sample smoke completes the
new pipeline but obtains sequence-macro MiD $1868.3186$; it is integration
evidence, not evidence of accuracy. The current bottleneck is therefore the
learned representation---a compatible dense JEPA pretrainer and RGB-event fusion
are still missing---rather than storage or basic pipeline execution.
\end{abstract}

\section{Scope and claim boundary}

The system targets signed continuous TTC from asynchronous events and later an
RGB-event multimodal extension. EvTTC~\cite{evttc2024} supplies local supervised
evaluation, whereas Garl-TTC/eAP~\cite{eapgarl2026} defines the principal external
comparison. Event-only and RGB-E are separate protocols.

A state-of-the-art claim requires matching modality, split, metric, and benchmark;
three predeclared seeds; a frozen checkpoint; label-free EvTTC prediction; and an
official eAP/CodaBench result. None of these external gates is currently closed.

\section{Historical EvTTC evidence}

The historical EventTubelet baseline was reconstructed byte-for-byte on its
original cache. Table~\ref{tab:history} separates that anchor from the later
matched architecture matrix.

\begin{table}[t]
\centering
\caption{Audited historical and grouped-CV evidence.}
\label{tab:history}
\small
\begin{tabular}{lrrr}
\toprule
Model/protocol & Score & RTE & MAE (s) \\
\midrule
BASE historical & -- & 8.1554\% & 0.3229 \\
A0 grouped CV & 0.58452 & 30.25\% & 1.011 \\
A1 Dense grouped CV & 0.59312 & 30.55\% & 1.007 \\
\bottomrule
\end{tabular}
\end{table}

Grouped-CV values are means across five folds and three seeds; their RTE standard
deviations are $0.52$ and $0.06$ percentage points, respectively. A1 improves
aggregate MAE by only $0.41\%$, worsens score and RTE, and costs approximately
$2.16\times$ inference latency. Attention Residuals, Object-KDA inspired by
delta-memory designs~\cite{kimik32026}, and bbox-ROI pooling also fail their
sequential gates. These negative results prevent complexity from being reintroduced
without a new predeclared hypothesis.

\section{Data and leakage controls}

EvTTC-32 uses indivisible sequence groups; no random window split is allowed. The
sealed EvTTC benchmark remains unopened. The local eAP inventory covers 40 of 46
expected sequences. A signed pilot uses 9/3 train/validation sequences and the full
local split uses 32/8, selected without EvTTC metrics.

The Garl join, timestamps, calibration, signed TTC metrics, and balanced sampling
were audited. TTC, depth, 3-D height, category, and masks are supervision or audit
metadata, never event-encoder inputs. CARLA was removed from the active inventory
after negative cross-domain transfer; compact negative summaries remain tracked.

\section{High-resolution model}

The active event-only path is
\begin{equation}
\begin{aligned}
X &\rightarrow \text{patches} \rightarrow \text{window attention}\\
  &\rightarrow \text{temporal mixer} \rightarrow \hat{\ttc}.
\end{aligned}
\end{equation}
Padding and validity masks preserve non-divisible borders. Optional $2\times2$
space-to-depth merging reduces token count before the block-causal temporal mixer,
which operates on $[B,T,P,D]$. Fixed learned queries pool retained patches for the
TTC head. A memory guard estimates attention allocations before execution. A BF16
query-pooling dtype defect found during the audit is now regression-tested.

The target is signed. Training minimizes Smooth L1 after
\begin{equation}
f(x)=\operatorname{sign}(x)\log(1+|x|),
\end{equation}
so separating objects are not incorrectly mapped to positive TTC. Model selection
uses sequence-macro MiD rather than a window-weighted average.

\section{JEPA status}

The intended objective predicts future dense target-encoder tokens at multiple
horizons, with stop-gradient and EMA. It must not use TTC or future labels. The
existing pooled legacy pretrainer is architecture-incompatible with the active
high-resolution token grid, so its alias now fails explicitly. The current
high-resolution candidate consequently starts from random initialization unless a
strictly compatible checkpoint is supplied.

This missing matched pretraining path is the principal scientific limitation. A
latent loss decrease in another architecture would not establish downstream TTC
transfer.

\section{Semantic capacity audit}

The real 192-dimensional SSL artifact has effective ranks 2.255, 1.095, and
5.105 for context, predictor, and target. The predictor is nearly one-dimensional,
but rank cannot identify whether that axis represents TTC, expansion, or sequence
identity. Healthy variance is not a semantic certificate, consistent with the
known slow-feature failure of predictive embeddings~\cite{slowfeatures2022}.

We therefore added a dataset-free falsifier. Five representation objectives use
identical observations and seeds 7/13/23; TTC and shortcut labels are withheld
during representation training and used only by probes fitted after freezing.
Table~\ref{tab:semantic} reports the fixed-per-sequence 12-bit shortcut.

\begin{table}[t]
\centering
\caption{Synthetic semantic-shortcut audit, three-seed mean.}
\label{tab:semantic}
\scriptsize
\begin{tabular}{lrrrr}
\toprule
Objective & Dyn. $R^2$ & MAE & Acc. & Dup. \\
\midrule
Repository variance & 0.15 & 0.39 & 0.84 & 1.93 \\
VISReg-style & 0.20 & 0.38 & 0.92 & 1.63 \\
Temporal residual & \textbf{0.72} & \textbf{0.29} & 0.65 & 1.06 \\
R$^2$ rate+dependence & 0.29 & 0.36 & 0.88 & 1.92 \\
Residual+R$^2$ & 0.48 & 0.34 & 0.68 & \textbf{0.68} \\
\bottomrule
\end{tabular}
\end{table}

Variance and the VISReg-style distributional control~\cite{visreg2026} retain the
slow shortcut. R$^2$-lite does not pass the predeclared TTC gate and is rejected
for production. Temporal residual prediction passes the slow-nuisance gate, but
in a frame-varying control its dynamic $R^2$ falls to $-0.05$, versus 0.74 for the
repository objective. It is therefore conditional rather than universal. The
minimal real candidate keeps a level channel for scale/content and adds a residual
channel for expansion/motion; it does not introduce rate, HSIC, or conditional-MI
estimators before a matched eAP screen.

\section{Cache-free training}

The trainer indexes raw HDF5 windows and materializes only the active batch. It
implements deterministic balanced sampling, BF16, gradient accumulation, clipping,
epoch-derived shuffle for reproducible resume, macro validation, and atomic best/
last checkpoints with provenance.

The screen uses dimension 32 and patch size 8. The full candidate uses dimension
192, patch size 16, batch 4 with accumulation 6, up to 30 epochs, and seeds
7/13/23. Full runs require a clean Git tree and share commit, configuration, data,
and split hashes before freezing.

\subsection{Storage audit}

A full dense cache was projected at approximately 455~GiB. A 256-sample shard was
materialized, SHA-verified, consumed, and removed. A 4096-sample attempt approached
11~GiB RAM without completing. The global-cache route is therefore excluded.

On 2 August, the rejected CARLA extraction (71.64~GiB), generated runs
(16.87~GiB), and feature caches (12.81~GiB) were deleted. Compact signed summaries
were retained. This recovered approximately 101~GiB in that cleanup pass.

\section{Current real smoke}

\begin{table}[t]
\centering
\caption{Raw event-only integration smoke; not promotable.}
\label{tab:smoke}
\small
\begin{tabular}{lr}
\toprule
Quantity & Value \\
\midrule
Train / validation samples & 16 / 16 \\
Epochs / seed & 1 / 7 \\
Elapsed time & 17.16 s \\
Sequence-macro MiD & 1868.3186 \\
Global MiD & 1941.2832 \\
Weighted RTE & 119.2892\% \\
Failure rate & 0\% \\
\bottomrule
\end{tabular}
\end{table}

The zero failure rate shows that the poor metric is not hidden exception filtering.
The run validates data flow, GPU execution, signed loss, metrics, and checkpoint
compatibility only. Scaling this formulation immediately would be expensive and
scientifically unjustified.

\section{Evaluation orchestration}

One runner separates the stages:
\begin{lstlisting}
train -> freeze -> evttc-predict
      -> evttc-score -> submission-validate
\end{lstlisting}
The full freeze selects among seeds only on Garl validation. EvTTC prediction
requires a hashed label-free manifest and rejects target-like fields before loading
the checkpoint or GPU. Scoring receives targets in a separate process. Submission
validation is offline and never counts as an official upload.

\section{Limitations and next gates}

Dense JEPA pretraining and RGB-E fusion remain unimplemented. The semantic audit
is synthetic and does not prove real eAP shortcut allocation. Six eAP sequences,
the real EvTTC Table-VI input manifest, full three-seed training, robustness,
calibration, final ONNX export, and a real demo are missing. Causal bbox-free
expansion/FoE has not surpassed A0, and ground-truth bbox/depth remain diagnostic
oracles. The system is not suitable for safety control.

The next economical sequence is: (1) overfit matched dense JEPA with level and
residual channels on 256 samples; (2) paired level-versus-level+residual screens
with frozen TTC, expansion, event-rate, and sequence-ID probes; (3) isolated RGB-E;
(4) real EvTTC predict/score; and only then (5) full 7/13/23 and official
eAP/CodaBench submission. These restrictions preclude a present SOTA claim.

\begin{thebibliography}{5}
\bibitem{evttc2024}
J.~Li et al.\ ``EvTTC: An Event Camera Dataset for Time-to-Collision
Estimation.'' arXiv:2412.05053, 2024.

\bibitem{eapgarl2026}
J.~Li et al.\ ``Toward Deep Representation Learning for Event-Enhanced Visual
Autonomous Perception: the eAP Dataset.'' arXiv:2603.16303, 2026.

\bibitem{kimik32026}
Kimi Team.\ ``Kimi K3: Open Frontier Intelligence.'' arXiv:2607.24653, 2026.

\bibitem{slowfeatures2022}
V.~Sobal et al.\ ``Joint Embedding Predictive Architectures Focus on Slow
Features.'' arXiv:2211.10831, 2022.

\bibitem{visreg2026}
``VISReg: Variance-Invariance-Sketching Regularization for JEPA Training.''
arXiv:2606.02572, 2026.
\end{thebibliography}

\end{document}
